\section{Interface utilisateur}
\label{sec:soft_ui}

Les couches \texttt{ui} et \texttt{screens} de la section~\ref{sec:archi} fournissent le cadre graphique générique et les écrans concrets. Cette section décrit ce cadre, soit la façon dont un écran est dessiné puis affiché, dont les écrans s'empilent et se cèdent la main, et dont les appuis de l'utilisateur y sont acheminés. L'interface est construite sans bibliothèque graphique, sur un tampon d'affichage et des primitives de dessin propres au projet.

\subsection{Rendu et tampon d'affichage}

Le rendu passe par un seul tampon d'affichage en mémoire, au format RGB565, soit seize bits par pixel. Pour les 176 sur 176 pixels de l'écran, il occupe 61\,952~octets. Toutes les primitives de dessin y écrivent, puis une opération d'envoi le transfère à l'écran. Une bibliothèque comme LVGL aurait fourni des composants tout faits, mais avec son empreinte mémoire et son style propres. Le projet a préféré un jeu de primitives réduit écrites pour cet écran uniquement, qui laissent la main sur le rendu pixel par pixel voulu pour l'interface.

Ce tampon est placé en \gls{psram}, selon l'arbitrage mémoire de la section~\ref{sec:archi}. L'envoi le lit ligne par ligne avec le processeur et recopie chaque ligne dans un petit tampon interne, seul ce dernier partant en transfert direct sur le bus SPI~\cite{ESP32_IDF_guide}. La mémoire externe coûte donc un accès processeur plus lent, mais jamais un transfert impossible.

Les primitives se limitent au nécessaire, soit le remplissage de rectangles, le tracé de contours, l'écriture de texte et le placement d'icônes. Chacune reste dans le cadre de l'écran et ignore les coordonnées hors champ, ce qui évite tout débordement. Le texte utilise une fonte matricielle de cinq sur sept pixels, dessinée dans une cellule de six sur huit et agrandie par un facteur entier, ce qui garde des contours nets à toutes les tailles. Les icônes sont définies sur un bit par pixel. Seuls les bits allumés sont peints, les autres laissent voir le fond, ce qui permet de les poser sur un contenu déjà dessiné.

L'écran est un module OLED à matrice passive, piloté par un contrôleur SSD1333 sur le bus SPI~\cite{NHD-1.91-176176UBC3}. Un pixel noir n'émet aucune lumière et ne consomme rien, propriété exploitée par le verrouillage de l'écran décrit à la section~\ref{sec:ui_verrouillage}.

\subsection{Écrans et navigation}

Tous les écrans exposent le même contrat, une table de quatre fonctions réunissant l'entrée, la sortie, le traitement d'un événement et le rendu. Chaque écran concret place cette table en tête de sa propre structure, si bien qu'un pointeur vers l'un vaut pointeur vers l'autre. Le cadre générique manipule ainsi tous les écrans de façon uniforme sans rien connaître de leur contenu.

Les écrans forment une pile de huit maximum, dont seul celui du dessus reçoit l'affichage et les entrées. Ouvrir un sous-menu ajoute un écran sur la pile, un retour le retire. Le cadre signale ces deux transitions à l'écran concerné, qui dispose ainsi d'un point d'accroche pour acquérir ses ressources et d'un autre pour les libérer.

Une barre de statut est enfin superposée à chaque écran. Le cadre la dessine en dernier, après le rendu de l'écran courant, de sorte qu'aucun écran ne peut l'effacer. Elle occupe la bande supérieure de l'image et y présente l'état du lien Bluetooth, la sortie audio active et le niveau de batterie. Ces indications restent qualitatives et se rafraîchissent au prochain appui plutôt que par un réveil propre, ce qui reste cohérent avec le fonctionnement événementiel de l'interface. Son apparence figure en annexe~\ref{ann:ecrans} avec les écrans concrets.

\subsection{Boucle d'interface}

La tâche d'interface exécute une boucle unique qui attend un événement, le remet à l'écran du dessus, puis présente l'image que celui-ci a dessinée. Cet ordre importe, le traitement de l'événement ayant pu changer l'écran courant. Dessiner sans relire la pile retarderait tout changement d'écran d'un appui. Au repos, cette attente est bloquante et ne consomme aucun cycle, la tâche étant retirée des tâches prêtes jusqu'au prochain appui, conformément au modèle de la section~\ref{sec:archi}.

La durée de cette attente n'est pas toujours infinie. Chaque écran déclare une période de rafraîchissement, nulle par défaut, qui borne l'attente. Une valeur nulle désigne un écran purement événementiel, redessiné au seul gré des appuis. Une valeur non nulle marque un écran vivant, comme les pages de statistiques, redessiné à intervalle régulier même sans appui. L'attente est de plus plafonnée à l'échéance de verrouillage, garantissant l'extinction de l'écran après le délai prévu en l'absence d'interaction.

Cette tâche initialise l'écran et le service d'entrée depuis son propre contexte, afin que leurs interruptions soient allouées sur le cœur~1 et n'empiètent pas sur le cœur~0 que sature la pile Bluetooth. Elle réarme par ailleurs le watchdog de tâche (\gls{twdt}) à chaque tour de boucle~\cite{ESP32_IDF_guide}. La tâche s'abonne au watchdog tant qu'un écran vivant est affiché. Celui-ci est redessiné plusieurs fois par seconde, la boucle réarme donc le watchdog bien avant son échéance. Si l'échéance est malgré tout atteinte, c'est que la tâche est réellement bloquée. Un écran statique ou verrouillé attend au contraire légitimement plusieurs dizaines de secondes un appui, durée qui serait prise à tort pour un blocage, la tâche s'en désabonne alors. Cet abonnement conditionnel évite que le fonctionnement événementiel voulu ne soit confondu avec une défaillance.

\subsection{Entrées}

Le service d'entrée ne pratique aucune scrutation. Sa tâche reste inactive tant qu'aucun bouton n'est actionné et ne consomme alors aucun cycle. Elle est relancée par une interruption, soit sur la ligne partagée de l'expanseur, soit sur la ligne propre du bouton A. Le gestionnaire d'interruption se borne à relancer la tâche, la lecture des boutons exigeant une transaction I\textsuperscript{2}C interdite en contexte d'interruption, comme l'a exposé la section~\ref{sec:archi}. La tâche relit les deux sources à chaque relance, ce qui la rend insensible au nombre d'interruptions reçues.

Les deux sources d'interruption tiennent au câblage. Les quatre boutons de navigation et le bouton de retour passent par l'expanseur, actifs à l'état bas, dont la ligne réagit à tout changement du port et se relâche à sa lecture~\cite{diodes_inc_PI4IOE5V9554}. Le bouton de validation dispose de sa propre broche, active à l'état haut. Son interruption est configurée sur les deux fronts, comme l'expanseur.

Les deux sources signalent ainsi les appuis comme les relâchements. La tâche a besoin des deux, faute de quoi un bouton maintenu puis relâché figerait l'antirebond sur un appui et ferait manquer le suivant. Elle lit à son réveil le port de l'expanseur, ce qui acquitte du même coup son interruption, puis l'état du bouton de validation, avant de confier ces relevés à l'antirebond.

L'antirebond est purement calculatoire et ne touche à aucun matériel. Il réunit les deux sources en un masque unique de boutons pressés, ne retient que les fronts d'appui au moyen d'un délai de garde qui absorbe les rebonds sans se déclencher sur un maintien ou un relâchement, et traduit chaque appui en un événement du vocabulaire commun, les quatre directions, la validation et le retour. Il tient enfin l'horodatage du dernier appui, dont se déduit la durée d'inactivité qui alimente le verrouillage de l'écran et la politique d'arrêt automatique de la section~\ref{sec:soft_power}.

\subsection{Économie d'affichage et verrouillage}
\label{sec:ui_verrouillage}

Passé un délai d'inactivité, l'écran est verrouillé. La boucle d'interface remplace alors l'écran courant par un écran de veille, une trame noire portant deux lignes d'indication en gris sombre. Ce verrouillage distingue un appareil en veille d'un appareil éteint. Son délai est une préférence de l'utilisateur configurée dans les réglages, une valeur nulle le désactivant entièrement.

Cette image de veille sert directement l'autonomie. L'écran étant émissif, un pixel noir n'émet aucune lumière et ne consomme rien, de sorte qu'une trame presque entièrement noire ne coûte que le courant de repos de la dalle et la poignée de pixels gris de l'indication. Seul le bouton de validation réveille l'écran et tout autre bouton est ignoré, ce qui neutralise les appuis accidentels en poche. La lecture audio et le réglage de volume, portés par d'autres tâches, restes fonctionnels.

Un dernier soin protège la dalle. Pendant le verrouillage, l'indication est redessinée périodiquement à une position légèrement décalée, selon un petit cycle de décalages, afin qu'aucun pixel ne vieillisse plus vite que ses voisins ni ne souffre de rémanence (burn-in)~\cite{viewsonic_oled_burnin}. Le coût de ces redraws est négligeable au regard de leur espacement.